\scanpage{68}
% archive sheet number: 67
% @ L1–9
de tantôt, pour \unsure{répérer|repérer} une op. de générisation $D'\longrightarrow D$ par un ``arc de \deleted{\unsure{spécial|spécialisation}}générisation'' sur $\partial\Delta_{D'}$ ; on prend maintenant, non la rive qui ``fait face'' à celle de $b$, mais par contre celle qui en résulte par transport parallèle. \deleted{Je \unsure{pensais}} \inserted{Quant à} l'isom. de transition $\underline{\omega}_{\Delta_D}\longrightarrow\underline{\omega}_{\Delta_{D'}}$ pour les orientations, faisant comme d'hab

% @ L10–11
\diagram{\begin{tikzcd}[ampersand replacement=\&,column sep=2em]
\omega\ \text{or. de }\Delta_D \arrow[r,leftrightarrow]\arrow[dr,to path={|- (\tikztotarget)\tikztonodes},"{\text{transport}}" near start,swap] \& \omega_D\ \text{or. de }D \& {}\\
{} \& \omega'\ \text{or. de }\Delta_{D'}\arrow[r,leftrightarrow] \& \omega_{D'}\ \text{or. de }D',
\end{tikzcd}}
% @ L12–14
on trouve à nouveau que $\omega_{D'}$ est \uline{opposée} à l'orientation \deleted{de $D$} déduite de $\omega_{D'}$ par transport parallèle !

% @ L15–end and R1–5
Pendant ces deux dernières semaines, c'est avec cette opération de vrillage que j'ai travaillé. Avant, j'avais pris \deleted{une} \uline{comme} iso. de transition pour les orientations le transport parallèle -- \deleted{\unsure{de façon}} et comme \deleted{\unsure{arcs}} vrillage, si je me rappelle bien, le vrillage $1^\circ)$ de $b$ avec $a'$ -- je croyais avoir interprété combinatoirement la construction de vrillement de $1^\circ)$, mais je m'étais mélangé les pinceaux de \uline{deux} façons, il me semble :

% @ R6–7
1) je m'étais trompé sur les isom. de transition pour les orientations, et

% @ R8–16
2) je m'étais trompé également sur l'identification correcte d'une op. de générisation \inserted{\unsure{pour|partt} $D'$} par un arc sur $\partial\Delta_{D'}$ -- en prenant au bonheur la chance celui qui \deleted{est} \deleted{\unsure{est}} se déduit de \uline{$b$} par transport parallèle, alors que c'est l'autre \deleted{qu'il} (l'antipodique, qu'il fallait prendre).

% @ R17–19
Pour résumer, les quatre possibilités sont indiqués dans la liste suivante
